\begin{tomtat}
\begin{itemize}
  \item Định nghĩa~\ref{def:ch07-ham-attribution} gọi \tn{hàm attribution}{attribution function}
  là phép gán điểm cho các phần tử của một tập đối tượng đã chọn. Ba nhánh chỉ
  khác nhau ở tập ấy: đặc trưng của đầu vào, điểm huấn luyện, hoặc thành phần
  bên trong mô hình.
  \item \tn{Attribution dữ liệu}{data attribution} và
  \tn{attribution thành phần}{component attribution} dùng lại đúng ba kỹ thuật
  của Phần~II là phép nhiễu, gradient và
  \tn{xấp xỉ tuyến tính}{linear approximation}. Ở nhánh dữ liệu, ba kỹ thuật ấy
  lần lượt là leave-one-out, \tn{hàm ảnh hưởng}{influence function} và
  Datamodel; ở nhánh thành phần là \tn{vá kích hoạt}{activation patching},
  attribution patching và COAR.
  \item Khung \tn{xấp xỉ hàm cục bộ}{local function approximation} trong phương
  trình~\ref{eq:ch07-lfa} đưa tám phương pháp attribution đặc trưng về một dạng
  duy nhất, khác nhau ở lân cận lấy mẫu và hàm mất mát. Việc mở rộng khung này
  sang hai nhánh còn lại mới là một giả thuyết.
  \item Khảo sát của Long và cộng sự xếp lĩnh vực theo phạm vi và theo thời điểm
  can thiệp. Không trục nào trong hai trục ghi lời giải thích có phản ánh đúng
  mô hình hay không.
  \item Bài báo quan điểm của Zhang và cộng sự ghi nhận ba nhánh dùng chung một
  bộ chỉ số đánh giá: fidelity,
  inverse fidelity và sparsity, tất cả đều dựa trên đánh giá phản thực và đều
  thiếu ground truth để đối chiếu. Nếu ghi nhận ấy đúng thì một sai sót trong bộ
  chỉ số cũng lan sang cả ba nhánh.
\end{itemize}
\end{tomtat}

\cauhoi
\begin{cauhoids}
  \item Hãy nêu tập đối tượng $E$ trong định nghĩa~\ref{def:ch07-ham-attribution}
  cho từng nhánh, rồi cho biết mỗi nhánh cần truy cập cái gì của quá trình huấn
  luyện: chỉ mô hình đã huấn luyện, hay cả tập dữ liệu và các checkpoint.
  \item Phương trình~\ref{eq:ch07-ham-anh-huong} thay việc huấn luyện lại bằng
  một khai triển Taylor. Khai triển ấy đòi những giả thiết nào, và vì sao học sâu
  thường không thỏa chúng?
  \item Trong phương trình~\ref{eq:ch07-lfa}, hai phương pháp có thể dùng chung
  lớp mô hình $\mathcal{G}$ và chỉ khác nhau ở $\mathcal{N}_x$. Hãy giải thích vì
  sao khác biệt đó đủ để hai phương pháp cho hai lời giải thích khác nhau cho
  cùng một dự đoán.
  \item Vá kích hoạt cần một lượt chạy với đầu vào đã làm hỏng. Hãy chọn hai
  trong năm cách làm hỏng mà chương nêu, rồi lập luận xem kết luận về một nơ-ron
  có thể đổi thế nào khi đổi cách làm hỏng.
  \item Đọc \enquote{Towards Unified Attribution in Explainable AI, Data-Centric
  AI, and Mechanistic Interpretability} của Zhang và cộng
  sự~\autocite{p18unified}. Phần bàn về các quan điểm đối lập nêu lý do nào để ba
  nhánh tiếp tục phát triển riêng, và các tác giả trả lời bằng loại lập luận nào?
  \item Đọc phần bàn về \tn{siêu suy luận}{meta-reasoning} trong khảo sát của
  Long và cộng sự~\autocite{p20trends}. Nếu một hệ thống giải thích chính quá trình suy luận
  của nó, ta cần bằng chứng gì để tin lời giải thích ấy? Hãy đối chiếu câu trả
  lời với các chỉ số fidelity ở cuối chương.
\end{cauhoids}
